% !TEX root = _main.tex

\section{楽曲データベース}
\label{app:database}
本研究で使用した楽曲データベースは，J-POPの楽曲を中心に収録したものである．
楽曲データベースの詳細を表\tabref{tab:database}に示す．


\begin{table}[tb]
	\caption{楽曲データベースの詳細}
	\label{tab:database}
	\centering
	\begin{tabular}{c|c|c|c|c|c|c|c|c|c} \hline
		\multicolumn{2}{|c|}{\multirow{2}{*}{ジャンル}} & \multicolumn{2}{c|}{\multirow{2}{*}{アーティスト}} & \multicolumn{2}{c|}{\multirow{2}{*}{楽曲数}} & \multicolumn{2}{c|}{\multirow{2}{*}{収録時間}} & \multicolumn{2}{c|}{\multirow{2}{*}{ファイル形式}} \\
		\multicolumn{2}{|c|}{}                      & \multicolumn{2}{c|}{}                        & \multicolumn{2}{c|}{}                     & \multicolumn{2}{c|}{}                      & \multicolumn{2}{c|}{}                        \\ \hline
		\multicolumn{2}{|c|}{J-POP}                 & \multicolumn{2}{c|}{Various Artists}         & \multicolumn{2}{c|}{100}                  & \multicolumn{2}{c|}{6:00:00}               & \multicolumn{2}{c|}{WAV}                     \\ \hline
	\end{tabular}
\end{table}